Many Tunisian student clubs and local associations rely on public social networks because they are easy to access and familiar to members. These tools are useful for visibility, but they are less suitable when the community needs internal coordination, clear roles, controlled membership, and accountable moderation. In this type of organization, the problem is not only to publish messages. The platform must also help members recognize each other, preserve trust, and allow responsible actors to explain administrative decisions after they are taken.

TunSociety was developed from this observation. The project proposes a web platform for organized communities that need both social interaction and governance. A standard member can create an account, manage a profile, publish posts, comment, react, search for members, send friend requests, use direct messages, and receive notifications. A moderator or administrator has access to additional workflows for reviewing content, following risk indicators, issuing warnings, freezing accounts, managing appeals, changing roles, and consulting audit logs. This separation gives the application a precise identity: TunSociety is an internal community-management platform, not a clone of a public social media website.

The solution was implemented as a full-stack application. The frontend uses Angular and TypeScript. The backend is built with ASP.NET Core and exposes REST APIs protected by JWT authentication. Data is stored in MySQL through Entity Framework Core. The project also integrates local AI-assisted moderation through Ollama and the Gemma model. This component evaluates submitted text, returns a risk-oriented assessment, and helps the backend choose whether content should be accepted, flagged for review, or blocked.

This report follows the progression of the project. The first chapter explains the context, the host environment, the project scope, the requirements, and the working method. The second chapter describes the conception phase through modules, actors, data structures, API contracts, access control, and sequence logic. The third chapter details the realization phase, including the development environment, implementation choices, validation scenarios, deployment preparation, maintenance, and future improvements. The report closes with a general conclusion and the consulted references.
